\section{2004秋}
\subsection{微分積分}
\prob{
  次の関数に対する偏導関数を求めなさい。
  \begin{enumerate}[label=(\arabic*), ref=(\arabic*), itemsep=5pt]
    \item $\disp f(x,y,z) = \cos^{-1}\bigl(x^3 y^2 z\bigr)$
    \item $\disp g(x,y) = x \tan^{-1}\biggl(\frac{y}{x}\biggr) - y \tan^{-1}\biggl(\frac{x}{y}\biggr)$
  \end{enumerate}
}

\prob{
  次の積分を求めなさい。
  \begin{gather}
    \int_{0}^{2\pi + i} z \sin z \,dz \qquad i = \sqrt{-1}
  \end{gather}
}

\prob{
  次の微分方程式を解きなさい。
  \begin{gather}
    y'' - 5y' + 6y = e^{-x}
  \end{gather}
}

\newpage
\subsection{線形代数}
\prob{
  式\eqref{eq:2004autumn_l1_1}に示した3次の実正方行列$\bA$が直交行列であることは、
  任意の3項実列ベクトル$\bx$に対して次の式を満足することと同値である。
  \begin{gather}
    \| \bA\bx \| = \| \bx \|
  \end{gather}
  ここに、$\|\,\|$はノルムである。$a,\,b,\,c,\,d$の値を一組求めよ。
  \begin{gather}
    \bA = \begin{pmatrix}
      a & \disp \frac{1}{\sqrt{2}} & \disp \frac{1}{\sqrt{3}} \\
      b & d & \disp \frac{1}{\sqrt{3}} \\
      c & e & \disp \frac{1}{\sqrt{3}}
    \end{pmatrix}\label{eq:2004autumn_l1_1}
  \end{gather}
}

\prob{
  次の連立一次方程式を解け。
  \begin{gather}
    6x - 2y + 2z = 2 \\
    2x + 6y + 4z = 4 \\
    4x - 6y - z = -1
  \end{gather}
}
